%----------------------------------------------------------------------------------------
% GONIOMETRISCHE FORMULES. Alle opgaven zijn eigen werk.
%
% Opbouw: Start, leerstof (theorie, werkwijze of uitleg, voorbeeld, korte oefening), en op een
% nieuwe pagina de Opdrachten (door elkaar, stijgende moeilijkheid), Koppel en Test jezelf.
% Labels van dit hoofdstuk beginnen met h2-.
%----------------------------------------------------------------------------------------

\mychapter{2}{Goniometrische formules}

\lkr{Leerplan Basis: LPD 11-13 en K2 (goniometrische functies), het deel over de som- en verschilformules en de verdubbelingsformules. Het bewijs van de verschilformule voor de cosinus is een mooi voorbeeld van een bewijs met een afstand: laat leerlingen bij elke regel zeggen waarom ze klopt. Sectie 3 (sinus en cosinus samenvoegen) bereidt de algemene sinusfunctie in hoofdstuk 3 voor.}

Is $\sin(\alpha + \beta)$ gelijk aan $\sin\alpha + \sin\beta$? Probeer $\alpha = \beta = \frac{\pi}{2}$: links staat $\sin\pi = 0$, rechts $1 + 1 = 2$. Nee dus. In dit hoofdstuk zoeken we wat $\sin(\alpha + \beta)$ en $\cos(\alpha + \beta)$ wel zijn. Daarmee bereken je exacte waarden voor nieuwe hoeken, en vereenvoudig je uitdrukkingen.

\begin{oefeningen}[Start]
\oef{h2-start-1}

\oef{h2-verken}
\end{oefeningen}

%========================================================================================
\section{Som- en verschilformules}

\begin{theorie}
\begin{stelling}{Som- en verschilformules}
\[\formule{\cos(\alpha - \beta) = \cos\alpha\cos\beta + \sin\alpha\sin\beta} \qquad \formule{\cos(\alpha + \beta) = \cos\alpha\cos\beta - \sin\alpha\sin\beta}\]
\[\formule{\sin(\alpha - \beta) = \sin\alpha\cos\beta - \cos\alpha\sin\beta} \qquad \formule{\sin(\alpha + \beta) = \sin\alpha\cos\beta + \cos\alpha\sin\beta}\]
\[\formule{\tan(\alpha \pm \beta) = \frac{\tan\alpha \pm \tan\beta}{1 \mp \tan\alpha\tan\beta}}\]
\end{stelling}
\begin{bewijs}
\naast{%
\textit{Verschil van de cosinus.} Neem de punten $P(\cos\alpha, \sin\alpha)$ en $Q(\cos\beta, \sin\beta)$ op de eenheidscirkel. Dan is
\bewijsregel{|PQ|^2 = (\cos\alpha - \cos\beta)^2 + (\sin\alpha - \sin\beta)^2}
en met de grondformule wordt dat
\bewijsregel{|PQ|^2 = 2 - 2(\cos\alpha\cos\beta + \sin\alpha\sin\beta).} Draai beide punten over $-\beta$: $Q$ gaat naar $Q'(1, 0)$ en $P$ naar $P'\bigl(\cos(\alpha - \beta), \sin(\alpha - \beta)\bigr)$. De afstand blijft dezelfde:
\bewijsregel{|P'Q'|^2 = \bigl(\cos(\alpha - \beta) - 1\bigr)^2 + \sin^2(\alpha - \beta)}
\bewijsregel{\phantom{|P'Q'|^2} = 2 - 2\cos(\alpha - \beta).}}{%
\begin{tikzpicture}[scale=1.3,font=\scriptsize]
  \eenheidscirkel
  \draw[accent] (0,0) -- (115:1) (0,0) -- (40:1);
  \draw[accent,thick] (115:1) -- (40:1);
  \fill[accent] (115:1) circle (.035) node[above left] {$P$};
  \fill[accent] (40:1) circle (.035) node[right] {$Q$};
  \draw[accent,dashed] (0,0) -- (75:1);
  \draw[accent,dashed,thick] (75:1) -- (1,0);
  \fill[accent] (75:1) circle (.035) node[above] {$P'$};
  \fill[accent] (1,0) circle (.035) node[below right] {$Q'$};
\end{tikzpicture}}
Gelijkstellen geeft de formule. De andere volgen eruit: vervang $\beta$ door $-\beta$ voor de som, en gebruik $\sin\gamma = \cos\bigl(\frac{\pi}{2} - \gamma\bigr)$ voor de sinus. De tangens is het quotiënt; deel teller en noemer door $\cos\alpha\cos\beta$.
\end{bewijs}
\end{theorie}

\voorbeeld{}{Bepaal exact $\cos\frac{\pi}{12}$.\label{vb:h2-exact}}{%
\uitleg{$\frac{\pi}{12}$ is geen bijzondere hoek, maar wel het verschil van twee: $\frac{\pi}{3} - \frac{\pi}{4} = \frac{4\pi - 3\pi}{12}$.}
\begin{stappenbreed}
\cos\tfrac{\pi}{12} = \cos\tfrac{\pi}{3}\cos\tfrac{\pi}{4} + \sin\tfrac{\pi}{3}\sin\tfrac{\pi}{4} & Verschilformule. \\
= \tfrac12 \cdot \tfrac{\sqrt2}{2} + \tfrac{\sqrt3}{2} \cdot \tfrac{\sqrt2}{2} = \tfrac{\sqrt2 + \sqrt6}{4} & Controle: $\approx 0{,}9659 \approx \cos 15^\circ$.
\end{stappenbreed}}

\voorbeeld{}{Gegeven: $\sin\alpha = \frac35$ en $\cos\beta = \frac{5}{13}$, met $\alpha$ en $\beta$ in het eerste kwadrant. Bereken $\sin(\alpha + \beta)$.\label{vb:h2-som}}{%
\uitleg{Zoek eerst de ontbrekende cosinus en sinus met de grondformule. In het eerste kwadrant zijn ze positief.}
\begin{stappenbreed}
\cos\alpha = \tfrac45,\ \sin\beta = \tfrac{12}{13} & Grondformule. \\
\sin(\alpha + \beta) = \tfrac35 \cdot \tfrac{5}{13} + \tfrac45 \cdot \tfrac{12}{13} = \tfrac{63}{65} & Somformule.
\end{stappenbreed}}

\begin{oefeningen}
\oef{h2-som}
\end{oefeningen}

%========================================================================================
\section{Verdubbelingsformules}

Neem $\beta = \alpha$ in de somformules. Dan krijg je formules voor de dubbele hoek.

\begin{theorie}
\begin{stelling}{Verdubbelingsformules}
\[\formule{\sin 2\alpha = 2\sin\alpha\cos\alpha} \qquad \formule{\tan 2\alpha = \frac{2\tan\alpha}{1 - \tan^2\alpha}}\]
\[\formule{\cos 2\alpha = \cos^2\alpha - \sin^2\alpha = 2\cos^2\alpha - 1 = 1 - 2\sin^2\alpha}\]
\end{stelling}
\begin{stelling}{Formules van Carnot}
\[\formule{\cos^2\alpha = \frac{1 + \cos 2\alpha}{2}} \qquad \formule{\sin^2\alpha = \frac{1 - \cos 2\alpha}{2}}\]
\end{stelling}
\end{theorie}

De drie vormen van $\cos 2\alpha$ volgen uit elkaar met de grondformule. De formules van Carnot zijn de laatste twee vormen, opgelost naar $\cos^2\alpha$ en $\sin^2\alpha$: ze ruilen een kwadraat in voor een dubbele hoek.

\voorbeeld{}{Gegeven: $\sin\alpha = \frac35$ en $\frac{\pi}{2} < \alpha < \pi$. Bereken $\sin 2\alpha$ en $\cos 2\alpha$.\label{vb:h2-dubbel}}{%
\uitleg{In het tweede kwadrant is $\cos\alpha = -\frac45$. Voor $\cos 2\alpha$ is de vorm met alleen $\sin\alpha$ het snelst.}
\begin{stappenbreed}
\sin 2\alpha = 2 \cdot \tfrac35 \cdot \bigl(-\tfrac45\bigr) = -\tfrac{24}{25} & \\
\cos 2\alpha = 1 - 2 \cdot \tfrac{9}{25} = \tfrac{7}{25} & Controle: $\bigl(\tfrac{24}{25}\bigr)^2 + \bigl(\tfrac{7}{25}\bigr)^2 = 1$.
\end{stappenbreed}}

\voorbeeld{}{Bewijs: $\dfrac{1 - \cos 2x}{\sin 2x} = \tan x$ (voor $\sin 2x \neq 0$).\label{vb:h2-identiteit}}{%
\uitleg{Werk het ingewikkeldste lid uit tot het andere. Kies voor $\cos 2x$ de vorm die de $1$ wegwerkt.}
\begin{stappenbreed}
\frac{1 - \cos 2x}{\sin 2x} = \frac{1 - (1 - 2\sin^2 x)}{2\sin x\cos x} & Verdubbelingsformules. \\
= \frac{2\sin^2 x}{2\sin x\cos x} = \frac{\sin x}{\cos x} = \tan x & Vereenvoudigen.
\end{stappenbreed}}

\begin{werkwijze}
Bij een identiteit: vertrek van één lid (meestal het ingewikkeldste) en werk het uit tot het andere. Schrijf alles met dezelfde hoek ($2x$ als $x$), en alles met $\sin$ en $\cos$. Kies de vorm van $\cos 2x$ die het meest vereenvoudigt.
\end{werkwijze}

\begin{oefeningen}
\oef{h2-dubbel}
\end{oefeningen}

%========================================================================================
\section{Sinus en cosinus samenvoegen}

Een som $a\sin x + b\cos x$ is altijd één verschoven sinus. Dat zie je met de somformule van achter naar voor.

\begin{theorie}
\begin{stelling}{Samenvoegen}
$\formule{a\sin x + b\cos x = R\sin(x + \varphi)}$ \quad met \quad $\formule{R = \sqrt{a^2 + b^2}}$, \quad $\formule{\cos\varphi = \frac{a}{R}}$ \ en \ $\formule{\sin\varphi = \frac{b}{R}}$.
\end{stelling}
\end{theorie}

Inderdaad: $R\sin(x + \varphi) = R\cos\varphi \sin x + R\sin\varphi \cos x = a\sin x + b\cos x$. Het punt $\bigl(\frac{a}{R}, \frac{b}{R}\bigr)$ ligt op de eenheidscirkel, dus zo'n hoek $\varphi$ bestaat.

\voorbeeld{}{Schrijf $\sin x + \sqrt3\cos x$ als $R\sin(x + \varphi)$. Wat is de grootste waarde?\label{vb:h2-samenvoegen}}{%
\uitleg{Hier is $a = 1$ en $b = \sqrt3$. Zoek de hoek met cosinus $\frac{a}{R}$ en sinus $\frac{b}{R}$.}
\begin{stappenbreed}
R = \sqrt{1 + 3} = 2 & \\
\cos\varphi = \tfrac12,\ \sin\varphi = \tfrac{\sqrt3}{2} \iff \varphi = \tfrac{\pi}{3} & Eerste kwadrant. \\
\sin x + \sqrt3\cos x = 2\sin\bigl(x + \tfrac{\pi}{3}\bigr) & De grootste waarde is $2$.
\end{stappenbreed}}

\begin{oefeningen}
\oef{h2-samenvoegen}
\end{oefeningen}

%========================================================================================
\opdrachtenkop

\begin{oefeningen}
\oef{h2-o1}

\oef{h2-o2}

\oef{h2-o3}

\oef{h2-o4}

\oef{h2-o5}

\oef{h2-o6}

\oef{h2-o7}

\oef{h2-o8}

\oef{h2-o9}
\end{oefeningen}

%----------------------------------------------------------------------------------------
% KOPPEL: verbanden met andere leerstof
%----------------------------------------------------------------------------------------
\needspace{20\baselineskip}
\opdrachtsectie{Koppel}

\lkr{Matrices: kijk na wat al gezien is. Het pakket legt de volgorde niet vast.}

\begin{oefeningen}

\oef{h2-koppel-3}
\end{oefeningen}

%----------------------------------------------------------------------------------------
\opdrachtsectie{Test jezelf}

\begin{oefeningen}
\oef{h2-t1}

\oef{h2-t2}

\oef{h2-t3}

\oef{h2-t4}

\oef{h2-t5}
\end{oefeningen}

\herhaalschema{%
  \fv{oef:h2-t1}{vb:h2-exact}{oef:h2-som}%
  \fv{oef:h2-t2}{vb:h2-som}{oef:h2-som}%
  \fv{oef:h2-t3}{vb:h2-dubbel}{oef:h2-dubbel}%
  \fv{oef:h2-t4}{vb:h2-identiteit}{oef:h2-dubbel}%
  \fv{oef:h2-t5}{vb:h2-samenvoegen}{oef:h2-samenvoegen}}

\ifzelfstudie
\section*{Maak je eigen samenvatting}
Sluit dit boek en schrijf in je schrift: de som- en verschilformules (met het idee van het bewijs), de verdubbelingsformules met de drie vormen van $\cos 2\alpha$, de formules van Carnot, en hoe je $a\sin x + b\cos x$ samenvoegt. Vergelijk daarna met de kaders.
\fi
